\section{Background and preliminaries}
\label{sec:preliminaries}

\subsection{The existential theory of the reals}

An existential sentence over the reals has the form
\begin{equation}
  \exists x_1,\ldots,x_m\in\R:\;
  \Psi(x_1,\ldots,x_m),
\end{equation}
where \(\Psi\) is a quantifier-free Boolean combination of polynomial sign
conditions with integer coefficients.  The decision problem ETR asks whether
such a sentence is true.  The class \(\ExistsR\) is the closure of ETR under
polynomial-time many-one reductions.  Standard containments give
\begin{equation}
  P\subseteq NP\subseteq\ExistsR\subseteq PSPACE.
\end{equation}
We use the discrete Turing model: coefficients and combinatorial structure are
encoded by finite binary strings, while quantified witnesses range over the
reals.  Surveys and algorithms for real feasibility may be found in
\cite{BasuPollackRoy2006,Renegar1992I,SchaeferCardinalMiltzow2024}; Canny's
algorithm supplies the PSPACE upper bound~\cite{Canny1988}.

\begin{definition}[Semialgebraic set]
A subset \(X\subseteq\R^d\) is semialgebraic if it is defined by a finite
Boolean combination of polynomial equalities and inequalities with real
coefficients.  A \emph{basic} semialgebraic set is defined by a conjunction of
such sign conditions.
\end{definition}

Existentially quantified square variables can replace weak nonnegativity:
\begin{equation}
  t\geq 0 \quad\Longleftrightarrow\quad \exists u\in\R:\;t=u^2.
  \label{eq:square-slack}
\end{equation}
This device permits an equality-only ambient variety whose real projection is
the intended basic semialgebraic set.  The projection operation remains part of
the semantics; the equality coordinate ring alone does not remember which
variables were introduced existentially.

\subsection{Coordinate rings and differentials}

Let \(k\) be a commutative base field and \(B\) a finitely generated
\(k\)-algebra.  The module of K\"ahler differentials \(\Om_{B/k}\) is the
\(B\)-module equipped with a universal \(k\)-derivation
\(d:B\to\Om_{B/k}\).  For a tower \(k\to A\to B\), there is a right-exact
transitivity sequence
\begin{equation}
  B\otimes_A\Om_{A/k}
  \longrightarrow \Om_{B/k}
  \longrightarrow \Om_{B/A}
  \longrightarrow 0.
  \label{eq:kahler-transitivity}
\end{equation}
These standard constructions are reviewed, for example, in
\cite{Eisenbud1995,StacksDifferentials}.

If
\begin{equation}
  B=k[x_1,\ldots,x_n]/(f_1,\ldots,f_m),
\end{equation}
then \(\Om_{B/k}\) admits the presentation
\begin{equation}
  B^m \xrightarrow{J_F^{\mathsf T}} B^n
  \longrightarrow \Om_{B/k}\longrightarrow 0,
  \label{eq:jacobian-presentation}
\end{equation}
where \(J_F=(\partial f_i/\partial x_j)\).  This is a module over the entire
quotient ring.  Evaluating at a real point \(p\) tensors the presentation with
the residue field at \(p\) and produces pointwise cotangent information.  The
symbolic module and any one evaluated vector space are distinct objects.

\subsection{Linear maps and intrinsic messages}

Let \(J_A:L_A\to X_A\) and \(J_{AS}:L_A\to X_S\) be linear maps over a field.
The intrinsic separator message is
\begin{equation}
  \cE(A\to S)
  :=J_{AS}(\kernel J_A)\subseteq X_S.
  \label{eq:intrinsic-linear-message}
\end{equation}
If the maps are represented by row blocks and \(P_A^\perp\) has rows spanning
the relevant left null space, then a matrix representing the same subspace is
\begin{equation}
  M_{A\to S}=P_A^\perp J_{AS}.
\end{equation}
The subspace, rather than the chosen basis matrix, is the invariant.

\subsection{Tree decompositions}

A tree decomposition of a graph \(G=(V,E)\) is a tree whose nodes carry bags
\(B_t\subseteq V\) such that every vertex appears in a connected family of
bags and every edge is contained in some bag.  Its width is
\(\max_t|B_t|-1\); the treewidth \(\tw(G)\) is the minimum width.  Bounded
treewidth often enables exact dynamic programming when each separator message
has a controlled representation~\cite{Bodlaender1996}.  Treewidth alone does
not bound the descriptive complexity of projected semialgebraic sets.

\subsection{Proof-status vocabulary}

To prevent a conditional dependency from being silently promoted, we use four
labels in this draft:
\begin{center}
\begin{tabular}{>{\bfseries}l p{0.72\linewidth}}
\toprule
Proved & A complete mathematical proof is present or a cited standard result
applies with its hypotheses checked.\\
Conditional & The implication is proved assuming explicitly named
certificates or lemmas.\\
Open & The statement is required by the target theorem but is not presently
derived.\\
Refuted & A counterexample disproves the statement as formulated.\\
\bottomrule
\end{tabular}
\end{center}
